% ICASSP 2027 manuscript: official SPS conference template (spconf.sty / IEEEbib.bst).
\documentclass{article}
\usepackage{spconf,amsmath,amssymb,booktabs,placeins,balance,hyperref,flushend}
\hypersetup{
  pdfauthor={Xianghui Meng and Jionghao Lin},
  pdftitle={When Does Reference Matching Transfer? Specification-Certified Audits of FIR/IIR Realizations},
  pdfsubject={ICASSP 2027 regular paper},
  pdfkeywords={Digital filters, FIR, IIR, specification testing, reference matching}
}

\title{WHEN DOES REFERENCE MATCHING TRANSFER? SPECIFICATION-CERTIFIED\\
AUDITS OF FIR/IIR REALIZATIONS}

\name{Xianghui Meng$^{1}$ and Jionghao Lin$^{1,2,\ast}$}
\address{$^{1}$The University of Hong Kong \\
$^{2}$Carnegie Mellon University \\
\texttt{margretmeng1020@gmail.com} \\
$^{\ast}$Corresponding author: \texttt{jionghao@hku.hk}}

\begin{document}
\maketitle

\begin{abstract}
Reference-based correctness scoring treats proximity to selected
realizations as a proxy for satisfaction of a digital-filter
specification, even when a magnitude mask admits many legitimate
FIR or IIR designs.
We audit that proxy on 20 magnitude-mask tasks (16 FIR, 4 IIR)
whose 412 constructed valid implementations are continuously
certified, using coefficient and magnitude-response catalogs of
observed valid realizations.
No observed valid coefficient reference recovers exact
single-center membership on any task; exact coefficient catalogs
require more than 10 references on all 20 tasks (median 23).
After those catalogs were frozen, a catalog-blind protocol produced
614 new continuously certified valid realizations.
The frozen coefficient catalogs accepted $66/614$ ($10.7\%$;
task-macro median $4.8\%$), whereas the frozen magnitude-response
catalogs accepted $585/614$ ($95.3\%$; task-macro median $100\%$).
Reference adequacy therefore depends on representation and on the
diversity of valid realizations that must be supported:
coefficient matching is tied to realization identity, while
magnitude-response matching is a substantially stronger---but not
universally exact---surrogate on this suite.
Code, frozen artifacts, and the reproduction entry point are at
\url{https://github.com/XianghuiMeng-1020/dsp-semantic-correctness}.
\end{abstract}

\begin{keywords}
Digital filters, FIR, IIR, specification testing, reference matching
\end{keywords}

\section{Introduction}
\label{sec:intro}

A digital-filter task is specified by passband and stopband
edges, ripple, attenuation, and, for IIR filters,
stability~\cite{oppenheim2010dsp,proakis2007dsp,rabiner1975theory,mitra2011dsp}:
a feasible set, not a unique coefficient vector.
Evaluation nevertheless often scores proximity to a golden
coefficient vector, a reference magnitude response, or a small
library catalog.
That surrogate asks whether a candidate is the \emph{same
realization}, not whether it satisfies the mask.
This paper audits whether such realization-reference oracles
transfer to alternative specification-valid FIR/IIR designs.

That Remez, window, least-squares, frequency-sampling, Butterworth,
Chebyshev, and elliptic constructions can meet one mask with
different taps or sections is classical filter
theory~\cite{kaiser1974window,harris1978windows,rabiner1970freqsamp,parks1972chebyshev,mcclellan1973computer,rabiner1975fir,herrmann1973linear,butterworth1930,constantinides1970,jackson1996filters,antoniou2006filters}
and is not our contribution.
The general test-oracle problem, specification-based conformance,
and prototype or set-cover selection are likewise
established~\cite{weyuker1982nontestable,barr2015oracle,huuhtanen2015dsp,hart1968cnn,chvatal1979setcover}.
We do not propose a new oracle theory or a new set-cover algorithm.
The question is whether coefficient or magnitude-response
references that exactly fit one specification-certified FIR/IIR
universe transfer to later catalog-blind implementations of the
same masks.

This paper makes three contributions.
(i)~We formulate correctness as membership in a
specification-defined valid set $\mathcal{V}_t$ and assemble a
continuously certified FIR/IIR evaluation universe.
(ii)~We quantify how finite catalogs of \emph{actual} valid
realizations recover that membership on the frozen universe,
including an exact observed-valid catalog-complexity diagnostic.
(iii)~On a catalog-blind prospective challenge of 614 independently
certified realizations, we show a strong
representation-dependent contrast: frozen coefficient catalogs
transfer poorly, frozen magnitude-response catalogs usually
transfer well, and admitting the new valids substantially expands
every exact coefficient catalog.

\section{Specification sets and reference catalogs}
\label{sec:form}

A task $t$ asks for an FIR impulse response or an IIR pair
$(b,a)$.
A predicate $S_t(h)\in\{0,1\}$ encodes the registered magnitude
mask and, for IIR designs, a strict pole-radius bound.
Correctness is membership
\begin{equation}
\mathcal{V}_t=\{h:S_t(h)=1\}.
\label{eq:vt}
\end{equation}
A reference catalog $\mathcal{R}$ is a finite set of
\emph{observed} specification-valid realizations.
With a distance $d$ and a common scalar threshold $\tau$, the
acceptance set is
\begin{equation}
\mathcal{A}_{\tau,\mathcal{R}}=\{h:d_{\mathcal{R}}(h)\le\tau\},
\quad
d_{\mathcal{R}}(h)=\min_{r\in\mathcal{R}}d(h,r).
\label{eq:A}
\end{equation}
The reference oracle declares ``correct'' iff
$h\in\mathcal{A}_{\tau,\mathcal{R}}$; the specification oracle
declares ``correct'' iff $h\in\mathcal{V}_t$.

All exact-recovery statements below are restricted to an explicit
finite universe $\mathcal{U}_t$.
Write $V=\mathcal{V}_t\cap\mathcal{U}_t$ and
$I=\mathcal{U}_t\setminus\mathcal{V}_t$, and define
\begin{equation}
D_V(\mathcal{R})=\max_{v\in V}\min_{r\in\mathcal{R}}d(v,r),
\quad
D_I(\mathcal{R})=\min_{i\in I}\min_{r\in\mathcal{R}}d(i,r).
\label{eq:DVDI}
\end{equation}
\emph{Observation~1.}
On a frozen finite universe, a common scalar threshold recovers
membership exactly if and only if
$D_V(\mathcal{R})<D_I(\mathcal{R})$.
The statement is evaluation geometry, not a theorem over all
filters.

The observed-valid \emph{reference catalog complexity}
$K^\star_{\mathrm{obs}}$ is the smallest
$\lvert\mathcal{R}\rvert$ of specification-valid occupants of
$V$ for which Observation~1 holds.
Computing that finite-universe diagnostic reduces to a standard
set-cover
optimization~\cite{chvatal1979setcover,hart1968cnn}; the
optimization itself is not our contribution.
$K^\star_{\mathrm{obs}}$ measures realizable reference burden, not
algorithmic novelty.

Two distances are used.
Coefficient distance is a relative $\ell_2$ after
canonicalization (FIR: trailing-zero and sign rules; IIR:
$a_0=1$, concatenated $(b,a)$).
Magnitude-response distance is the RMSE of $\lvert H\rvert$ on the
constrained pass and stop bands at $N_f=131072$ frequency
samples.
Centers must be actual observed valid realizations, not arbitrary
Euclidean points.

A length-31 Kaiser-window low-pass design and a length-45
weighted-Chebyshev (Parks--McClellan) design can both satisfy the
same passband-ripple and stopband-attenuation mask while differing
substantially in length, tap values, and phase-delay structure, so
their coefficient vectors sit far apart under $\ell_2$ even when
their constrained-band magnitude responses nearly coincide.
Coefficient distance is therefore sensitive to which valid
realization was observed; magnitude-response distance is
comparatively insensitive to that choice.
Intuitively, $K^\star_{\mathrm{obs}}$ grows when $V$ is
geometrically interleaved with $I$ under the chosen distance: a
single center under-covers $V$ or over-covers into $I$ whenever
some valid realization is farther from every candidate center than
some invalid realization is from its nearest center.
Because coefficient distance uses each realization's own tap or
section parameterization as the operative geometry, and different
design families place their occupants in different regions of that
space, a coefficient-space center separating one family's occupants
from invalids need not separate another family's occupants from the
same invalids---the source of the observed $K^\star{>}10$
requirement.
IIR tasks add a further partition: the strict pole-radius bound
constrains each design family's $(b,a)$ occupants to a
family-specific region of coefficient space before any mask
criterion is applied, so an IIR coefficient catalog inherits both
the mask-driven and the stability-driven separation the FIR case
does not.

\section{Certified universe and prospective protocol}
\label{sec:protocol}

The base suite comprises 20 non-unique magnitude-mask tasks:
FIR low-pass, high-pass, band-pass, and band-stop at
$f_s=8$ and $16\,\mathrm{kHz}$, each loose and tight ($16$ FIR);
and IIR low-pass and high-pass at $8\,\mathrm{kHz}$ ($4$ IIR).
The frozen constructed valid corpus contains $412$
implementations ($336$ FIR, $76$ IIR).

Final labels are not self-certified by the construction checker.
For FIR, $P(x)=\lvert H\rvert^2$ with $x=\cos\omega$ is certified
continuously over each constrained band by exact-rational
Bernstein or Sturm sign reasoning.
For IIR, we certify strict Schur stability (pole radius
$<0.999$) and the sign of $P_B(x)-C P_A(x)$ on every constrained
band.
All $412/412$ constructed valids received continuous
specification certification; there were no valid$\to$invalid
contradictions, and all frozen mechanism and boundary invalids
certified invalid.
Cosine endpoint handling is a numerical/rational certificate, not
a machine-checked continuous proof.
Implementation details live in the repository.

The eight prospective design families were realized with standard
SciPy signal-processing routines~\cite{virtanen2020scipy}:
\texttt{remez} (Parks--McClellan), \texttt{firls}, \texttt{firwin2}
(frequency sampling), and \texttt{firwin} (Kaiser/Hann windowing)
for FIR; \texttt{butter}, \texttt{cheby1}, \texttt{cheby2}, and
\texttt{ellip} for IIR.
No family was excluded a priori, and each family's outputs were
certified by the same continuous procedure as the base catalogs
before being counted toward the $614$ frozen valids.

Three questions organize the audit.
RQ1: when does realization-reference matching recover
specification membership on this frozen universe?
RQ2: how large a realizable catalog is required for exact
recovery?
RQ3: do exact base-fitted catalogs transfer to later
implementations that satisfy the same masks?

We first froze the base catalogs and their max-safe thresholds
(the largest $\tau$ that accepts no frozen base invalid).
We then protocol-locked generation across all 20 tasks and eight
standard families---Remez, FIRLS, frequency sampling, and
windowed FIR; Butterworth, Chebyshev~I,~II, and elliptic
IIR---scheduled $960$ attempts, continuously certified occupants,
removed exact coefficient duplicates, and froze $614$ new
nonduplicate valid realizations ($500$ FIR, $114$ IIR; every
task; at least $13$ per task) without accessing catalog
identities, reference distances, transfer scores, or ambient
centers.
Only then did we evaluate transfer.
An earlier same-order probe set participated in catalog fitting
and is not this holdout.
The $614$ occupants are prospectively generated, catalog-blind,
specification-certified realizations, not all possible filters.

\FloatBarrier
\section{Results}
\label{sec:results}
\begin{table*}[t]
\centering
\caption{Catalog-blind prospective transfer on $614$
specification-certified realizations, by magnitude-mask task
(L/T: loose/tight).
Coeff.\ and Resp.\ are the frozen coefficient- and
magnitude-response-catalog acceptance rates; $N$ is the number of
prospective valids for that task.}
\label{tab:transfer}
\setlength{\tabcolsep}{4pt}
\footnotesize
\begin{tabular}{@{}lrrr@{}|lrrr@{}}
\toprule
Task & $N$ & Coeff.\ & Resp.\ & Task & $N$ & Coeff.\ & Resp.\ \\
\midrule
FIR LP, L, 8k  & 44 & $6.8\%$  & $100\%$ & FIR BS, L, 8k  & 40 & $32.5\%$ & $97.5\%$ \\
FIR LP, L, 16k & 44 & $4.5\%$  & $100\%$ & FIR BS, L, 16k & 40 & $5.0\%$  & $97.5\%$ \\
FIR LP, T, 8k  & 26 & $3.8\%$  & $100\%$ & FIR BS, T, 8k  & 13 & $0.0\%$  & $100\%$ \\
FIR LP, T, 16k & 26 & $0.0\%$  & $100\%$ & FIR BS, T, 16k & 13 & $0.0\%$  & $100\%$ \\
FIR HP, L, 8k  & 41 & $19.5\%$ & $97.6\%$ & IIR LP, L, 8k  & 41 & $17.1\%$ & $78.0\%$ \\
FIR HP, L, 16k & 41 & $2.4\%$  & $97.6\%$ & IIR LP, T, 8k  & 18 & $0.0\%$  & $100\%$ \\
FIR HP, T, 8k  & 24 & $50.0\%$ & $91.7\%$ & IIR HP, L, 8k  & 42 & $4.8\%$  & $100\%$ \\
FIR HP, T, 16k & 24 & $12.5\%$ & $87.5\%$ & IIR HP, T, 8k  & 13 & $0.0\%$  & $100\%$ \\
FIR BP, L, 8k  & 42 & $2.4\%$  & $100\%$ &                &    &          & \\
FIR BP, L, 16k & 42 & $4.8\%$  & $100\%$ &                &    &          & \\
FIR BP, T, 8k  & 20 & $15.0\%$ & $70.0\%$ &                &    &          & \\
FIR BP, T, 16k & 20 & $30.0\%$ & $75.0\%$ &                &    &          & \\
\bottomrule
\end{tabular}
\end{table*}

\subsection{Base-universe adequacy}
On the frozen $412$-valid universe, no single observed valid
coefficient reference recovers exact membership (20/20
canonical and 20/20 best-observed non-separable).
A single magnitude-response reference fails on $19/20$ canonical
and $18/20$ best-observed tasks (Table~\ref{tab:base}).
Exact observed-valid coefficient catalogs exist on $20/20$ tasks,
but $K^\star>10$ on every task (median $23$).
Library-only catalogs recover membership on $0/20$ tasks in both
representations.
Response $K^\star$ has median $17.5$ and equals $1$ on two tasks.

An arbitrary Euclidean center separates 19/20 coefficient and
20/20 response base universes (one IIR low-pass task has no
coefficient center), so the observed failures should not be
read as intrinsic one-center inseparability; the restriction
to actual realizations is consequential.

\begin{table}[t]
\caption{Exact recovery on the frozen base universe.
``Non-sep.'' is a finite-universe single-center failure.
$K^\star$ is the minimum observed-valid catalog size.}
\label{tab:base}
\centering
\setlength{\tabcolsep}{4pt}
\begin{tabular}{@{}lcccc@{}}
\toprule
Repr. & Can.\ non-sep. & Best obs. & med.\ $K^\star$ & Lib.\ exact \\
\midrule
Coeff. & $20/20$ & $20/20$ & $23$ & $0/20$ \\
Resp.  & $19/20$ & $18/20$ & $17.5$ & $0/20$ \\
\bottomrule
\end{tabular}
\end{table}

\subsection{Prospective transfer}
This is the primary empirical result (Table~\ref{tab:transfer},
Table~\ref{tab:xfer}).
Frozen coefficient $K^\star_{\mathrm{obs}}$ catalogs accepted
$66/614$ prospective valids (pooled $10.7\%$; task-macro mean
$10.6\%$, median $4.8\%$; range $0$--$50\%$).
Every task transferred below $75\%$ ($20/20$); none reached
$95\%$.
FIR and IIR task-macro means were $11.8\%$ and $5.5\%$.
The same verdict is obtained under the predeclared midpoint
threshold, so the coefficient failure is not an artifact of
max-safe $\tau$.

Frozen magnitude-response catalogs accepted $585/614$ (pooled
$95.3\%$; task-macro mean $94.6\%$, median $100\%$).
Fifteen tasks reached at least $95\%$; one tight FIR band-pass
task transferred at $70\%$.
FIR and IIR macros were both about $94.6\%$.
Response transfer is generally strong and still task-specific.

Coefficient transfer failed across all eight generator families
(family rates $3.6\%$--$18.2\%$); response transfer for those
families ranged from $85.7\%$ to $100\%$.
Coefficient transfer reached $0\%$ on exactly three tasks---both
tight FIR band-stop tasks and the tight FIR low-pass task at
$16\,\mathrm{kHz}$---where the certified prospective valids sit far,
in tap space, from every frozen base-catalog member.
Response transfer was weakest on the two tight FIR band-pass tasks
($70\%$ at $8\,\mathrm{kHz}$, the only task below $75\%$; $75\%$ at
$16\,\mathrm{kHz}$), followed by the tight FIR high-pass tasks
($87.5\%$--$91.7\%$) and the loose IIR low-pass task ($78.0\%$).

Enlarging the coefficient catalog on the \emph{base} universe does
not rescue prospective transfer.
Canonical $K{=}1$, published $K{\in}\{3,5\}$, and all-library
coefficient oracles accepted $0/614$; the best observed single
coefficient reference accepted $7.5\%$; the exact
$K^\star_{\mathrm{obs}}$ catalog accepted $10.7\%$.
The corresponding response figures are $38.9\%$ (canonical) and
$95.3\%$ ($K^\star_{\mathrm{obs}}$).
Increasing coefficient-catalog size enough to fit the original
finite universe therefore does not produce strong prospective
coefficient transfer.

\begin{table}[t]
\caption{Prospective transfer of frozen base catalogs to $614$
catalog-blind certified-valid realizations.}
\label{tab:xfer}
\centering
\setlength{\tabcolsep}{3.6pt}
\begin{tabular}{lccccc}
\toprule
Repr. & Accepted & Pooled & Macro med. & ${\ge}95\%$ & ${<}75\%$ \\
\midrule
Coeff. & $66/614$ & $10.7\%$ & $4.8\%$ & $0/20$ & $20/20$ \\
Resp.  & $585/614$ & $95.3\%$ & $100\%$ & $15/20$ & $1/20$ \\
\bottomrule
\end{tabular}
\end{table}

\subsection{Catalog maintenance}
Admitting the $614$ prospective valids and recomputing exact
coefficient catalogs raises the median $K^\star$ from $23$ to
$55$ (median relative growth $52.5\%$).
Every task has $M^\star>0$: every minimum-size expanded
coefficient catalog requires at least one newly admitted
realization (median $M^\star=15.5$, maximum $37$).
In-sample catalog exactness therefore does not eliminate later
support-expansion burden.
Response catalogs grow much less (median $\Delta K=0$;
$M^\star>0$ on $9/20$ tasks).

\section{Discussion}
\label{sec:disc}

Coefficient matching retains realization-specific
parameterization: distinct Remez, window, or elliptic
occupants of the same mask can sit far apart in tap or
$(b,a)$ space, so a catalog exact on one sample can miss later
valid designs ($66/614$).
Magnitude-response matching quotients out much of that
coefficient-level variability and aligns more directly with the
pass/stop mask $S_t$, which is why the same catalogs accept
$585/614$ occupants.
The contrast is a signal-processing representation fact, not
merely that one rate is larger.
Response matching remains a surrogate: it is not $100\%$, one
task transferred below $75\%$, and it does not prove mask
satisfaction.

A golden coefficient vector, or a finite coefficient catalog
fitted to one implementation corpus, should not be read as a
correctness oracle for a non-unique filter-design task without an
adequacy audit.
Where magnitude-response matching transfers, it can be a useful
surrogate.
The criterion used to define correctness in this study is still
$S_t$.
For evaluation practice, a reference-based filter-design checker
should record which representation its catalog uses and should
periodically recompute $K^\star$-style adequacy diagnostics as new
valid realizations are observed, rather than treating one-time
exactness on a fixed corpus as a permanent guarantee: no coefficient
catalog earns that guarantee here, and even a favorable
magnitude-response catalog needs periodic reappraisal.

This audit targets magnitude-mask FIR/IIR filters realized as
direct-form coefficient vectors; cascaded, lattice, or multirate
realizations~\cite{vaidyanathan1993multirate} were not part of the
certified universe and are left to future evaluation-methodology
work.
The study is finite: 20 magnitude-mask tasks, not general DSP
software, and eight standard designer families, not all
implementations.
A secondary catalog-blind invalid challenge, covering $12/20$ tasks
($310$ realizations) and therefore not a balanced 20-task holdout,
found response false acceptance ($96/310$, $31.0\%$) exceeding
coefficient false acceptance ($38/310$, $12.3\%$) on that subset;
full per-task results are in the repository.
No infinite-universe claim is made, and magnitude-response
matching is not asserted to be sound or complete in general.

\section{Conclusion}
\label{sec:conc}

Reference matching should be audited as a surrogate for a
digital-filter specification, not assumed equivalent to it.
Across $614$ prospectively generated valid FIR/IIR realizations,
coefficient catalogs that were exact on one implementation
universe transferred poorly, while magnitude-response catalogs
were substantially more stable.
Correctness evaluation should distinguish specification
compliance from realization identity and choose the reference
representation accordingly.

Returning to the three questions: RQ1 (does realization-reference
matching recover specification membership on the frozen universe?)
is answered negatively for a single coefficient reference and only
partially for a single response reference; RQ2 (how large a
realizable catalog is required?) needs more than ten coefficient
references on every task, with a median of $23$; RQ3 (does an
exact base-fitted catalog transfer?) is answered asymmetrically,
with coefficient catalogs generalizing poorly ($66/614$) and
magnitude-response catalogs generalizing well but imperfectly
($585/614$).
The same audit methodology---freeze catalogs, generate a
catalog-blind challenge, and measure transfer by representation---
could be applied to other specification types, such as
phase-response or group-delay masks, without requiring a new
correctness definition.

Code and frozen artifacts are in the accompanying public
repository, where the certified base and prospective corpora, the
$K^\star$/RCC diagnostics, and a one-line reproduction entry point
are frozen and tagged so that every number in
Section~\ref{sec:results} can be regenerated independently.

\clearpage
\noindent\textbf{Compliance with ethical standards.}
No funding was received for conducting this study.
The authors have no relevant financial or nonfinancial interests
to disclose.
The study did not involve human participants, animal subjects, or
personally identifiable data.

\bibliographystyle{IEEEbib}
\bibliography{refs}

\end{document}
